\documentclass[12pt]{exam}

% --- Packages ---
\usepackage[utf8]{inputenc}
\usepackage[T1]{fontenc}
\usepackage[french]{babel}
\usepackage{amsmath, amssymb}
\usepackage{geometry}
\usepackage{tikz}
\usepackage[most]{tcolorbox}
\usepackage{lmodern}
\usepackage{pgfplots}
\pgfplotsset{compat=1.18}
\usepackage{mdframed}
\usepackage{bm}
\usepackage{xcolor}
\usepackage{sectsty}
\usetikzlibrary{calc}

\usepackage[utf8]{inputenc}
\usepackage[T1]{fontenc}
\usepackage[dvipsnames]{xcolor}
\usepackage{tkz-tab}
\usepackage{multido}

\geometry{margin=2cm}
\setlength{\parindent}{0pt}
\footer{}{\thepage}{\it{G.M.}}

% =====================
% Document
% =====================
\begin{document}

\begin{center}
    {\Large \textbf{Exercices demi groupe}}\\[2mm]
    {\large 2nde}\\[2mm]
    \textit{Fonctions affines et vecteurs}
\end{center}

\textbf{Exercice : Fonctions affines}

On considère la fonction $f$ dont la représentation graphique est donnée ci-dessous.

\begin{center}
\begin{tikzpicture}[scale=0.8]

% Grille
\draw[step=1cm, gray!40, very thin] (-2,-4) grid (6,5);

% Axes
\draw[->] (-2.2,0) -- (6.2,0) node[right] {$x$};
\draw[->] (0,-4.2) -- (0,5.2) node[above] {$y$};

% Graduations
\foreach \x in {-2,-1,1,2,3,4,5}
    \draw (\x,0.1) -- (\x,-0.1) node[below] {\small $\x$};

\foreach \y in {-3,-2,-1,1,2,3,4}
    \draw (0.1,\y) -- (-0.1,\y) node[left] {\small $\y$};

% Fonction f
\draw[thick, smooth] (-2,-1) -- (6,3);
\node at (4.6,3.3) {$\mathcal{C}_f$};

\end{tikzpicture}
\end{center}

\begin{questions}

\question Refaire la figure sur le cahier.

% -----------------------------
\question Trouver l'expression de $f(x)$.

% -----------------------------
\question Résoudre graphiquement \textbf{ET} algébriquement l’inéquation :
\[
f(x) \geq 1
\]

% -----------------------------
\question On considère la fonction $g$ définie par :
\[
g(x) = -2x + 5
\]

\begin{parts}
\part Tracer la représentation graphique de la fonction $g$ sur le graphique ci-dessus.
\part Déterminer graphiquement les coordonnées du point d’intersection des droites représentant $f$ et $g$.
\end{parts}

% -----------------------------
\question Résoudre graphiquement \textbf{ET} algébriquement l’inéquation :
\[
f(x) \leq g(x)
\]

\end{questions}


\end{document}